\documentclass[12pt,a4paper]{article}
\usepackage[utf8]{inputenc}
\usepackage[spanish]{babel}
\usepackage{geometry}
\usepackage{listings}
\usepackage{xcolor}
\usepackage{parskip}
\usepackage{amsmath}
\usepackage{hyperref}

\geometry{margin=2.5cm}

\lstset{
  basicstyle=\ttfamily\small,
  keywordstyle=\color{blue}\bfseries,
  commentstyle=\color{gray},
  breaklines=true,
  frame=single,
  language=Java,
  numbers=left,
  numberstyle=\tiny\color{gray},
  literate={á}{{\'a}}1 {é}{{\'e}}1 {í}{{\'\i}}1 {ó}{{\'o}}1 {ú}{{\'u}}1 {ñ}{{\~n}}1 {Á}{{\'A}}1 {É}{{\'E}}1 {Ó}{{\'O}}1
}

\title{\textbf{Conecta 4 con Negamax}\\
  \large Documentación Técnica}
\author{Inteligencia Artificial}
\date{}

\begin{document}
\maketitle
\tableofcontents
\newpage

\section{Descripción del Proyecto}
Implementación del juego \textbf{Conecta 4} con inteligencia artificial basada en \textbf{Negamax}. El jugador humano compite contra la CPU en un tablero de 6 filas por 7 columnas (configurable). El objetivo es alinear 4 fichas consecutivas en horizontal, vertical o diagonal antes que el oponente.

\section{Reglas del Juego}

\begin{itemize}
  \item El tablero tiene 6 filas × 7 columnas (configurable hasta 10×10).
  \item Los jugadores alternan turnos colocando fichas en columnas.
  \item Las fichas caen por gravedad hasta la fila más baja disponible.
  \item Gana el primero en alinear \textbf{4 fichas} consecutivas en cualquier dirección.
  \item Si el tablero se llena sin ganador, el juego termina en empate.
\end{itemize}

\section{Representación del Estado}

El tablero se almacena como una \textbf{matriz bidimensional} \texttt{int[filas][columnas]} donde:
\begin{itemize}
  \item \texttt{0} = celda vacía
  \item \texttt{1} = ficha del jugador
  \item \texttt{2} = ficha de la CPU
\end{itemize}

La gravedad se modela buscando la fila más baja vacía de cada columna:
\begin{lstlisting}
int filaDisponible(int col) {
    for (int f = filas-1; f >= 0; f--)
        if (tablero[f][col] == 0) return f;
    return -1; // columna llena
}
\end{lstlisting}

\section{Algoritmo Negamax}

\subsection{Función de evaluación}
Negamax requiere evaluar tableros no terminales. La función \texttt{evaluar()} asigna puntuación analizando \textbf{ventanas de 4 celdas} en todas las direcciones:

\begin{itemize}
  \item 4 fichas propias: \texttt{+1000}
  \item 3 fichas propias + 1 vacía: \texttt{+50}
  \item 2 fichas propias + 2 vacías: \texttt{+10}
  \item 3 fichas del oponente + 1 vacía: \texttt{-80}
\end{itemize}

\subsection{Pseudocódigo Negamax con poda alfa-beta}
\begin{lstlisting}
función negamax(tablero, profundidad, alfa, beta, jugador):
    si profundidad == 0 o juegoTerminado(tablero):
        retornar evaluar(tablero, jugador)

    mejorPuntuacion = -infinito
    para cada columna disponible (0..6):
        fila = filaDisponible(tablero, col)
        tablero[fila][col] = jugador
        puntuacion = -negamax(tablero, profundidad-1,
                              -beta, -alfa, oponente)
        tablero[fila][col] = 0
        mejorPuntuacion = max(mejorPuntuacion, puntuacion)
        alfa = max(alfa, puntuacion)
        si alfa >= beta: break  // poda
    retornar mejorPuntuacion
\end{lstlisting}

\subsection{Poda alfa-beta}
La poda alfa-beta elimina ramas del árbol que no pueden influir en el resultado final, reduciendo la complejidad de $O(b^d)$ a $O(b^{d/2})$ en el mejor caso.

\section{Verificación de Victoria}

Se comprueban las 4 direcciones posibles para cada celda del tablero:

\begin{lstlisting}
int[][] direcciones = {{0,1},{1,0},{1,1},{1,-1}};
// horizontal, vertical, diagonal /, diagonal \

boolean hay4Iguales(int fila, int col, int jugador) {
    for (int[] dir : direcciones) {
        int count = 0;
        for (int k = 0; k < 4; k++) {
            int f = fila + k*dir[0];
            int c = col  + k*dir[1];
            if (dentroDelTablero(f,c) &&
                tablero[f][c] == jugador) count++;
        }
        if (count == 4) return true;
    }
    return false;
}
\end{lstlisting}

\section{Interfaz Gráfica}

La GUI fue construida con Java Swing:
\begin{itemize}
  \item \textbf{Panel de tablero}: dibujado con \texttt{paintComponent()} usando círculos de colores.
  \item \textbf{Detección de columna}: al hacer clic, se calcula la columna por la posición X del mouse.
  \item \textbf{Animación}: la ficha se dibuja inmediatamente al colocarla con \texttt{repaint()}.
  \item \textbf{Selección de color}: diálogo inicial para que el jugador elija rojo o amarillo.
\end{itemize}

\section{Flujo de Ejecución}

\begin{enumerate}
  \item El jugador selecciona color y nivel de dificultad.
  \item El jugador hace clic en una columna; la ficha cae a la fila más baja.
  \item Se verifica victoria o empate.
  \item La CPU ejecuta Negamax y coloca su ficha en la columna óptima.
  \item Se repite hasta estado terminal.
\end{enumerate}

\end{document}
